\setcounter{section}{17}
\section{Rozwiązanie fundamentalne laplasjanu i całki osobliwe}
\label{sec:dt-singular}
\providecommand{\atlasref}[2]{#2}
Por. \atlasref{PDE-07::thm:pde7-laplace}{RRC-07: rozwiązanie fundamentalne laplasjanu}
oraz \atlasref{PDE-16::lem:pde16-distr-kernel}{RRC-16: jądra osobliwe}.
\subsection{Rozwiązanie fundamentalne operatora Laplace'a}
\begin{twierdzenie}[Rozwiązanie fundamentalne operatora $-\Delta$]
\label{thm:dt-laplace}
Niech $d\ge1$ i niech $\omega_d$ oznacza pole powierzchni
sfery jednostkowej $S^{d-1}\subset\R^d$ dla $d\ge2$.
Funkcja lokalnie całkowalna
\begin{equation}\label{eq:dt-laplace-fundamental}
 \Phi_d(x)=\begin{cases}
 -|x|/2,&d=1,\\
 (2\pi)^{-1}\log(1/|x|),&d=2,\\
 ((d-2)\omega_d)^{-1}|x|^{2-d},&d\ge3,
 \end{cases}
\end{equation}
spełnia $-\Delta T_{\Phi_d}=\delta_0$ na $\R^d$.
\end{twierdzenie}
\begin{proof}
Wartość $\Phi_d(0)$ w wymiarach $d\ge2$ ustalamy
dowolnie, na przykład równą $0$.
W wymiarze jeden $\Phi_1$ jest lokalnie całkowalna
i jej pierwsza pochodna dystrybucyjna jest funkcją
$-\tfrac12\operatorname{sgn}x$, ponieważ funkcja jest ciągła
i kawałkami liniowa. Skok tej pochodnej w zerze wynosi $-1$.
Z twierdzenia~\ref{thm:jump20} wynika zatem
$\partial_x^2T_{\Phi_1}=-\delta_0$.

Dla $d\ge3$ całkowanie w zmiennych biegunowych daje
dla $r>0$
\[
 \int_{B_r(0)}|\Phi_d(x)|\,\mathrm d\lambda_d(x)
 =\frac1{d-2}\int_0^r s\,\mathrm d\lambda_1(s)
 =\frac{r^2}{2(d-2)}.
\]
Dla $d=2$ mamy $\omega_2=2\pi$ i dla $0<r<1$
\[
 \begin{aligned}
 \int_{B_r(0)}|\Phi_2(x)|\,\mathrm d\lambda_2(x)
 &=\int_0^r s\log(1/s)\,\mathrm d\lambda_1(s)\\
 &=\frac{r^2}{2}\log(1/r)+\frac{r^2}{4}.
 \end{aligned}
\]
Granica dolnego wyrazu pierwotnej jest zerowa,
gdyż $s^2\log(1/s)\to0$ dla $s\downarrow0$.
W obu przypadkach całka dąży do zera przy $r\downarrow0$.
Z gładkości poza zerem wynika więc lokalna
całkowalność na całej $\R^d$.

Niech $F(r)$ oznacza profil radialny $\Phi_d$ dla $r>0$.
Bezpośrednie różniczkowanie obu wzorów daje
\[
 F'(r)=-\frac1{\omega_dr^{d-1}},\qquad
 F''(r)=\frac{d-1}{\omega_dr^d}.
\]
Dla $x\ne0$, $r=|x|$,
$\partial_jr=x_j/r$ oraz
$\partial_j^2r=1/r-x_j^2/r^3$.
Z reguły łańcuchowej otrzymujemy
\[
 \partial_j^2\Phi_d(x)
 =F''(r)\frac{x_j^2}{r^2}
 +F'(r)\left(\frac1r-\frac{x_j^2}{r^3}\right).
\]
Sumując po $j$ i korzystając z $\sum_j x_j^2=r^2$,
dostajemy
\[
 \Delta\Phi_d(x)
 =F''(r)+\frac{d-1}{r}F'(r)=0.
\]
Dla $\varphi\in\mathcal D(\R^d)$ wybieramy $R>0$
tak, aby $\supp\varphi\Subset B_R(0)$.
Dla $0<\varepsilon<R$ stosujemy drugą tożsamość
Greena (wynikającą z twierdzenia o dywergencji
dla pola $\Phi_d\nabla\varphi-\varphi\nabla\Phi_d$) na pierścieniu
$A_{\varepsilon,R}=B_R(0)\setminus\overline{B_\varepsilon(0)}$
(do części rzeczywistej i urojonej testu osobno).
Ponieważ $\Delta\Phi_d=0$ na pierścieniu,
otrzymujemy
\[
 -\int_{A_{\varepsilon,R}}\Phi_d\Delta\varphi
 \,\mathrm d\lambda_d
 =\int_{\partial A_{\varepsilon,R}}
 (\varphi\partial_\nu\Phi_d-\Phi_d\partial_\nu\varphi)
 \,\mathrm d\mathcal H^{d-1}.
\]
Na zewnętrznej sferze znikają test i jego pochodne.
Na wewnętrznej sferze $\nu=-x/|x|$, zatem
$\partial_\nu\Phi_d=-F'(\varepsilon)
=1/(\omega_d\varepsilon^{d-1})$.
W konsekwencji
\[
 -\int_{A_{\varepsilon,R}}\Phi_d(x)\Delta\varphi(x)
 \,\mathrm d\lambda_d(x)
 =\int_{\partial B_\varepsilon}\varphi(x)
 \frac{1}{\omega_d\varepsilon^{d-1}}\,\mathrm d\mathcal H^{d-1}(x)
 -\int_{\partial B_\varepsilon}\Phi_d(x)
 \partial_\nu\varphi(x)\,\mathrm d\mathcal H^{d-1}(x).
\]
Miara powierzchniowa sfery o promieniu $\varepsilon$
wynosi $\omega_d\varepsilon^{d-1}$. Po odjęciu
$\varphi(0)$ od pierwszej całki moduł różnicy
jest więc co najwyżej
$\sup_{|x|=\varepsilon}|\varphi(x)-\varphi(0)|$,
co dąży do zera z ciągłości testu.
Moduł drugiej całki nie przekracza
\[
 \|\nabla\varphi\|_\infty\,
 |F(\varepsilon)|\,\omega_d\varepsilon^{d-1}.
\]
Dla $d\ge3$ iloczyn bez normy gradientu wynosi
$\varepsilon/(d-2)$, dla $d=2$ i $\varepsilon<1$
wynosi $\varepsilon\log(1/\varepsilon)$.
Obie wielkości dążą do zera.
Ponieważ sfery mają miarę Lebesgue'a zero
i test znika poza $B_R(0)$, w całce po $\R^d$
pozostaje uwzględnić tylko wkład z $B_\varepsilon(0)$.
Jego moduł jest ograniczony przez
$\|\Delta\varphi\|_\infty
\int_{B_\varepsilon(0)}|\Phi_d(x)|\,\mathrm d\lambda_d(x)$,
co dąży do zera z wcześniejszych obliczeń.
Ze znaku $(-1)^2=1$ przy drugiej pochodnej
dystrybucyjnej otrzymujemy
\[
 \langle-\Delta T_{\Phi_d},\varphi\rangle
 =-\int_{\R^d}\Phi_d(x)\Delta\varphi(x)
 \,\mathrm d\lambda_d(x)=\varphi(0).
\]
Jest to żądana równość.
\end{proof}

\subsection{Jądra drugich pochodnych potencjału Newtona}
Niech $\omega_{d-1}=\mathcal H^{d-1}(S^{d-1})$ oznacza pole
jednostkowej sfery, a $\Gamma$ rozwiązanie fundamentalne $-\Delta$:
$-\Delta\Gamma=\delta_0$. Z twierdzenia~\ref{thm:dt-laplace} dla $d\ge3$ można wziąć
$\Gamma(x)=\bigl((d-2)\omega_{d-1}\bigr)^{-1}|x|^{2-d}$,
a dla $d=2$ funkcję $\Gamma(x)=-(2\pi)^{-1}\log|x|$.
Dla $x\ne0$ w każdym z obu przypadków
$\partial_i\Gamma(x)=-x_i/(\omega_{d-1}|x|^d)$.
Różniczkując względem $x_j$, otrzymujemy
\begin{equation}\label{eq:dt-cz-kernel}
 K_{ij}(x):=\partial_{ij}\Gamma(x)
 =\frac{d x_ix_j-\delta_{ij}|x|^2}
 {\omega_{d-1}|x|^{d+2}},\qquad 1\le i,j\le d,
\end{equation}
gdzie $\delta_{ij}$ oznacza deltę Kroneckera.
Jądro jest gładkie poza zerem i spełnia
$K_{ij}(tx)=t^{-d}K_{ij}(x)$ dla $t>0$.
Różniczkowanie tej tożsamości daje jednorodność gradientu
stopnia $-d-1$. Z ograniczoności obu funkcji na zwartej
sferze wynika $|K_{ij}(x)|\le C_d|x|^{-d}$ oraz
$|\nabla K_{ij}(x)|\le C_d|x|^{-d-1}$.
Dla $i\ne j$ odbicie względem $\theta_i=0$ zmienia znak
$\theta_i\theta_j$ i zachowuje miarę powierzchniową.
Dla $i=j$ całki z $\theta_i^2$ są równe przez symetrię
permutacji współrzędnych, a ich suma wynosi $\omega_{d-1}$.
Stąd
\[
 \int_{S^{d-1}}\theta_i\theta_j\,\mathrm d\mathcal H^{d-1}
 =\frac{\omega_{d-1}}d\delta_{ij},\qquad
 \int_{S^{d-1}}K_{ij}(\theta)\,\mathrm d\mathcal H^{d-1}(\theta)=0.
\]
Każde jądro jest niezerowe na pewnym otwartym podzbiorze
sfery (dla $i=j$ weźmy $\theta=e_i$, dla $i\ne j$
$\theta=(e_i+e_j)/\sqrt2$). Współrzędne biegunowe dają
\[
 \int_{B_1(0)}|K_{ij}(z)|\,\mathrm d\lambda_d(z)
 =\int_{S^{d-1}}|K_{ij}(\theta)|\,\mathrm d\mathcal H^{d-1}(\theta)
 \int_0^1 r^{-1}\,\mathrm d\lambda_1(r)=+\infty .
\]
Jeżeli $f$ jest ciągła i $f(x)\ne0$, to $|f(x-z)|$
jest oddzielone od zera w otoczeniu zera. Zwykła całka
Lebesgue'a z $K_{ij}(z)f(x-z)$ nie jest więc bezwzględnie
zbieżna, nawet dla takich $f\in C_c^\infty$.

\subsection{Główna wartość całki i dystrybucja związana z jądrem}

\begin{definicja}[Główna wartość całki]
\label{def:dt-principal-value}
Niech $f\in C_c^1(\R^d)$, a $K$ będzie jądrem klasy $C^1$ poza
zerem, jednorodnym stopnia $-d$ i o zerowej średniej sferycznej.
Definiujemy
\[
 T_Kf(x)=\operatorname{p.v.}\!\int_{\R^d}K(x-y)f(y)\,\mathrm d\lambda_d(y)
 :=\lim_{\varepsilon\downarrow0}
 \int_{|x-y|>\varepsilon}K(x-y)f(y)\,\mathrm d\lambda_d(y).
\]
Dla $0<\varepsilon<1$ podstawiamy $z=x-y$.
Jednorodność i zerowa średnia dają zerową całkę z $K$
po każdym pierścieniu. Zatem po odjęciu $f(x)$
w pierścieniu $\varepsilon<|z|<1$ dostajemy w granicy
\[
 \begin{aligned}
 T_Kf(x)&=\int_{|z|<1}K(z)(f(x-z)-f(x))\,\mathrm d\lambda_d(z)\\
 &\quad+\int_{|z|\ge1}K(z)f(x-z)\,\mathrm d\lambda_d(z).
 \end{aligned}
\]
Obie całki są bezwzględnie zbieżne: w pierwszej
$|f(x-z)-f(x)|\le\|\nabla f\|_\infty |z|$,
więc całka radialna jest ograniczona przez stałą razy
$\int_0^1\mathrm d\lambda_1(r)$; druga ma moduł
nie większy niż $C\|f\|_1$. Uzasadnia to istnienie
granicy dla każdego $x$ oraz
$|T_Kf(x)|\le C(\|\nabla f\|_\infty+\|f\|_1)$.
\end{definicja}

Dystrybucję $\operatorname{p.v.}K$ określamy przez
\[
 \begin{aligned}
 \langle\operatorname{p.v.}K,\varphi\rangle
 &=\int_{|z|<1}K(z)(\varphi(z)-\varphi(0))\,\mathrm d\lambda_d(z)\\
 &\quad+\int_{|z|\ge1}K(z)\varphi(z)\,\mathrm d\lambda_d(z).
 \end{aligned}
\]
Dla testów o nośniku w ustalonym zwartym $L$ moduł
działania jest ograniczony przez
$C\|\nabla\varphi\|_\infty+C\lambda_d(L)\|\varphi\|_\infty$.
Liniowość i oszacowanie przez półnormy na każdym zwartym
nośniku dowodzą ciągłości dystrybucyjnej na mocy
twierdzenia \ref{thm:distribution-estimate19}.
Jest to także dystrybucja temperowana. Dla
$\varphi\in\mathcal S$ pierwsza całka jest ograniczona
przez stałą razy $\|\nabla\varphi\|_\infty$,
a druga przez
$C Q_{d+1,0}(\varphi)\int_1^\infty r^{-1}(1+r)^{-d-1}\,\mathrm d\lambda_1(r)$,
z dodatkowym stałym czynnikiem powierzchni sfery.
To daje ciągłość na $\mathcal S$.

\begin{lemat}[Hessian rozwiązania fundamentalnego jako dystrybucja]
\label{lem:dt-cz-distr-kernel}
Dla każdego $1\le i,j\le d$ zachodzi w $\mathcal D'(\R^d)$
\begin{equation}\label{eq:dt-cz-distr-kernel}
 \partial_{ij}\Gamma
 =\operatorname{p.v.}K_{ij}-\frac{\delta_{ij}}d\delta_0.
\end{equation}
W szczególności dla $f\in C_c^\infty(\R^d)$ i $u=\Gamma*f$,
\begin{equation}\label{eq:dt-cz-repr}
 \partial_{ij}u=T_{ij}f-\frac{\delta_{ij}}d f,
 \qquad T_{ij}=T_{K_{ij}}.
\end{equation}
\end{lemat}
\begin{proof}
Niech $\varphi\in C_c^\infty(\R^d)$, $R>1$ tak duże,
aby nośnik $\varphi$ leżał w $B_R(0)$, oraz
$D_\varepsilon=\{x:\varepsilon<|x|<R\}$.
Z lokalnej całkowalności $\Gamma$ działanie
$\partial_{ij}\Gamma$ jest granicą całki
$\int_{D_\varepsilon}\Gamma\partial_{ij}\varphi\,\mathrm d\lambda_d$.
Całkując przez części najpierw względem $x_i$, potem
$x_j$, otrzymujemy
\[
 \begin{aligned}
 \int_{D_\varepsilon}\Gamma\partial_{ij}\varphi\,\mathrm d\lambda_d
 &=\int_{D_\varepsilon}K_{ij}\varphi\,\mathrm d\lambda_d\\
 &\quad+\int_{\partial D_\varepsilon}
 \Gamma n_i\partial_j\varphi\,\mathrm d\mathcal H^{d-1}\\
 &\quad-\int_{\partial D_\varepsilon}
 (\partial_i\Gamma)n_j\varphi\,\mathrm d\mathcal H^{d-1}.
 \end{aligned}
\]
Na zewnętrznej sferze oba wyrazy brzegowe znikają.
Na wewnętrznej sferze $x=\varepsilon\theta$ mamy
$n=-\theta$, $\partial_i\Gamma(x)=
-\theta_i/(\omega_{d-1}\varepsilon^{d-1})$,
a miara powierzchniowa skaluje się przez
$\varepsilon^{d-1}$. Ostatni wyraz wynosi więc
\[
 -\frac1{\omega_{d-1}}\int_{S^{d-1}}\theta_i\theta_j
 \varphi(\varepsilon\theta)\,\mathrm d\mathcal H^{d-1}(\theta)
 \longrightarrow-\frac{\delta_{ij}}d\varphi(0).
\]
Przedostatni wyraz jest ograniczony co do modułu przez
$C\|\partial_j\varphi\|_\infty\varepsilon|\log\varepsilon|$
dla $d=2$ i przez $C_d\|\partial_j\varphi\|_\infty\varepsilon$
dla $d\ge3$, więc znika w granicy.
Pierwsza całka dąży do działania $\operatorname{p.v.}K_{ij}$.
Daje to~\eqref{eq:dt-cz-distr-kernel}.
Z twierdzeń~\ref{thm:smoothing21} i
\ref{thm:convolution-derivative21} splot
$u=\Gamma*f$ jest gładki, $-\Delta u=\delta_0*f=f$ oraz
$\partial_{ij}u=(\partial_{ij}\Gamma)*f$.
Działanie $\operatorname{p.v.}K_{ij}$ na teście
$z\mapsto f(x-z)$ jest dokładnie $T_{ij}f(x)$,
a $\delta_0*f=f$. Otrzymujemy~\eqref{eq:dt-cz-repr}.
\end{proof}

\subsection{Mnożniki związane z jądrem Newtona}
\begin{lemat}[Oszacowanie w $L^2$]
\label{lem:dt-cz-l2}
Dla $f\in C_c^\infty(\R^d)$ zachodzi
$\|T_{ij}f\|_2\le C_d\|f\|_2$. Operator $T_{ij}$
przedłuża się jednoznacznie do ograniczonego operatora na
$L^2(\R^d)$.
\end{lemat}
\begin{proof}
Najpierw wykażemy przynależność do $L^2$, bez jednostajnego
oszacowania normy przez $\|f\|_2$. Z definicji głównej
wartości $T_{ij}f$ jest ograniczona na całej przestrzeni.
Niech $\operatorname{supp}f\subset B_R(0)$, $R\ge1$.
Dla $|x|>2R$ mamy $|x-y|\ge|x|/2$ na nośniku $f$,
więc całka nie jest osobliwa i
$|T_{ij}f(x)|\le C_d|x|^{-d}\|f\|_1$.
Całka kwadratu tego ograniczenia poza $B_{2R}$ jest
skończona, gdyż $\int_{2R}^\infty r^{-d-1}
\,\mathrm d\lambda_1(r)<\infty$.
Na $B_{2R}$ wystarcza ograniczoność. Stąd $T_{ij}f\in L^2$,
a z~\eqref{eq:dt-cz-repr} także $H_{ij}:=\partial_{ij}u\in L^2$.

Z $-\Delta u=f$ i przemienności pochodnych mamy
$-\Delta H_{ij}=\partial_{ij}f$ dystrybucyjnie.
Obie strony są dystrybucjami temperowanymi; równość
w $\mathcal S'$ wynika z gęstości $C_c^\infty$ w $\mathcal S$
(twierdzenie~\ref{thm:density-schwartz17}). Twierdzenie~\ref{thm:fourier-auto-tempered22} daje
\[
 |\xi|^2\widehat H_{ij}=-\xi_i\xi_j\widehat f
 \qquad\text{w }\mathcal S'(\R^d).
\]
Obie strony są funkcjami lokalnie całkowalnymi, ponieważ
$\widehat H_{ij},\widehat f\in L^2$ z twierdzenia
Plancherela. Z jednoznaczności dystrybucji regularnych
(stwierdzenie~\ref{prop:regular-injective19}) mamy równość
prawie wszędzie. Dzieląc dla $\xi\ne0$, otrzymujemy
$\widehat H_{ij}(\xi)=-(\xi_i\xi_j/|\xi|^2)\widehat f(\xi)$
prawie wszędzie. Nie ma dodatkowego składnika skupionego
w zerze, bo $\widehat H_{ij}$ jest funkcją $L^2$.
Z~\eqref{eq:dt-cz-repr} wynika
$\widehat{T_{ij}f}=m_{ij}\widehat f$, gdzie
$m_{ij}(\xi)=\delta_{ij}/d-\xi_i\xi_j/|\xi|^2$
poza zerem, a wartość w zerze jest nieistotna.
Mamy $|m_{ij}|\le1+1/d\le2$. Z twierdzenia
~\ref{thm:plancherel22} przy przyjętej konwencji
\[
 \|T_{ij}f\|_2=\|m_{ij}\widehat f\|_2
 \le 2\|\widehat f\|_2=2\|f\|_2.
\]
Argument działa również dla $d=2$, bez wymagania, aby
sam potencjał $u$ należał do $L^2$. Liniowość, gęstość
$C_c^\infty$ w $L^2$ i zupełność $L^2$ dają jednoznaczne
ciągłe przedłużenie (gęstość w $L^2$ jest wykorzystana w dowodzie
twierdzenia~\ref{thm:plancherel22}).
\end{proof}
